\documentclass[10pt]{article}
\usepackage[margin=20mm]{geometry}
\usepackage{amsmath,amssymb,fontspec,unicode-math,xurl,hyperref}
\IfFontExistsTF{DejaVu Serif}{\setmainfont{DejaVu Serif}}{\setmainfont{Latin Modern Roman}}
\setmathfont{latinmodern-math.otf}
\setlength{\parindent}{0pt}\setlength{\parskip}{6pt}\sloppy
\begin{document}
\section*{Scientific verification report: I25-40168 (0060)}
Supported impact: Minor, for demonstrated source-characterization and printed-prompt defects. Central unresolved item: Needs Evidence for the calculation of prediction-based regret. Author adoption has not been requested.\par
\subsection*{1. Identity, version and scope}
Matthew Riemer, Zahra Ashktorab, Djallel Bouneffouf, Payel Das, Miao Liu, Justin D. Weisz and Murray Campbell, Position: Theory of Mind Benchmarks are Broken for Large Language Models, ICML 2025, PMLR 267:82091-82130 [E1-E2]. The primary reviewed object is the 40-page publisher-linked PDF. Page numbers below are one-based physical PDF pages, matching printed article pages.\par
Primary PDF SHA-256:
`98f89707b86dbb0569b4dd1a0834deea35618daf1f812098777c37c9075a6354`\par
The original 8,779,996 bytes match the source-acquisition record and remain preserved. The review covered the main argument, Definitions 2.1-2.2, the regret comparator, Tables 1-10, conceptual Appendices A/B, experimental descriptions in Appendix C, selected prompt images and selected learning curves. Figure captions throughout Appendix C were read; all curves were not numerically reconstructed.\par
The arXiv record identifies v1 on 27 December 2024, v2 on 5 February 2025, v3 on 6 June 2025 and v4 on 12 June 2025. A separate v4 PDF and its TeX-source archive were preserved. Targeted comparison confirms the same selected nonzero IPD values and comparator description; this is not a complete version diff. The TeX was read statically and was not executed [E3-E4].\par
\subsection*{2. Outcome and evidentiary standard}
The paper advocates evaluating actions and adaptation in addition to explicit predictions about other agents. It defines literal and functional theory of mind operationally, permits model-free implementations, and recognizes that prediction alone can suffice for insight-only applications. Its position is therefore strongest when interpreted as a warning against granting agency solely on the strength of predictive benchmarks.\par
Two bounded defects are demonstrated: an inaccurate description of cited interactive studies, and a prompt example that mixes prediction of the partner’s action with selection of the agent’s own action. Both support Minor corrections. Their historical runtime effects have not been established.\par
A separate, consequential question concerns `Delta\_ToM / T`. A stagewise interpretation contradicts the reported single-action IPD values, but the general policy notation permits a different, counterfactual planning interpretation. The available description does not identify the actual estimator sufficiently to choose between them. This remains Needs Evidence, with an exact conditional diagnostic and an explicit counterexample to overclaiming recorded below.\par
A demonstrated invalid core metric or implementation would support Major even if the broad recommendation survived or repair were easy. The present limitation is identification of the actual computation, not a requirement to prove that corrected experiments worsen performance. No empirical replication, author-code execution or model/API call was performed.\par
\subsection*{3. Finding 1: cited studies do contain interaction}
Judgment: Minor, literature-characterization error.\par
Claim and location. Section 3, p. 5, groups persona simulation, trust behavior and macroeconomic simulation studies, then says these tasks all lack interactivity. The matching primary sources contain interaction, including feedback used in subsequent decisions.\par
Cited study | Preserved version and locator | Directly relevant evidence\par
Xie et al., trust behavior | arXiv:2402.04559v1, PDF p. 3, Game 6; Appendix C, PDF p. 21 | A repeated trust game uses the same players. The next-round prompt supplies the previous amounts sent, received and returned before eliciting another decision.\par
Ha et al., CloChat | arXiv:2402.15265v1, Section 5.3, PDF p. 9 | Participants conduct live conversations with CloChat and ChatGPT; the procedure permits natural dialogue without imposed time or length limits.\par
Li et al., macroeconomic agents | arXiv:2310.10436v1, Section 4.2, PDF pp. 7-8 | A memory pool retains prior environments and decisions; periodic reflection informs later decisions.\par
Xie is the strongest witness because the interaction involves the same partner, outcome feedback and another action. Ha corroborates human-agent interaction; Li corroborates sequential environment interaction. These are evidence for one correction, not three independently inflated findings [E5-E7].\par
Strongest alternative reading. The authors may mean that these studies do not isolate reward-optimal partner adaptation or compare explicit prediction quality with the quality of resulting actions. That would be a narrower and potentially useful distinction. It does not make the no-interaction description accurate. The present checks do not establish that these earlier studies already implement the proposed paired metrics or achieve low functional regret.\par
Repair and impact. Acknowledge existing interactive work and specify which controlled diagnostic remains missing. Explain the value added by paired prediction/action measurements and regret. This changes the motivation and literature characterization; it does not establish that the paper’s observed prediction/action discrepancies are false.\par
The preserved Ha v1 has different presentation and pagination from the final CHI paper, but its title, DOI and procedure identify the study. Li v1 is the earlier four-author work corresponding to the cited 2023 title; the later EconAgent revision is not silently substituted. Xie v1 already contains the repeated-game protocol. All three versions predate the target.\par
\subsection*{4. Finding 2: Figure 13 mixes two prediction targets}
Judgment: Minor for the printed example; historical experimental impact Needs Evidence.\par
Section 4.2, p. 7, cites Figure 13 as the CoT 3-Shot variant of the action-selection prompt. In Figure 13, p. 23, the three examples reason about the focal agent’s own best response. Yet the final prompt states that the focal agent has already chosen Option J and asks which option the other player chose. It then asks for the best action. The final target is therefore ambiguous between predicting the partner and choosing the focal agent’s action.\par
This is visible in the original page image, so it is not a text-extraction artifact. Figures 4-7 separately illustrate action selection and literal prediction; the distinction is central to the experiment. Figure 13 does not clearly specify the action-selection procedure cited in the main text.\par
Strongest alternative reading. The figure could be a presentation mistake, or could show the prediction variant without an adequate label. Either possibility limits the demonstrated defect to the printed protocol. No execution trace or experiment implementation was located that establishes which template generated the reported CoT 3-Shot rows.\par
Repair. Supply separately labeled templates for choosing one’s own action and predicting the partner’s action. For an action-selection template, remove the statement that the action is already chosen and ask directly which action the agent should take. Check the actual run configuration before deciding whether only the figure needs correction or the affected results need recalculation. The proposed textual repair has not been applied to the source paper or validated in an LLM experiment.\par
\subsection*{5. Central pending item: reconcile prediction-based IPD regret}
Status: Needs Evidence; conditional contradiction established, unconditional implementation defect not established.\par
\subsubsection*{What must be reconciled}
Sections 2.1-2.2, pp. 3-4, describe a rational policy based on predicted partner behavior and ground-truth payoffs, with its regret reported as `Delta\_ToM / T`. The active TeX text gives the same description at lines 280 and 284. Table 2, p. 7, nevertheless reports nonzero IPD comparator regrets against single-action partners, including 0.098 for LM prompting and 1.342 for QA prompting. Table 4, p. 9, reports 0.360 for the DeepSeek single-action IPD row; Tables 8-9 also contain nonzero values.\par
Table 7, p. 18, gives the following focal-agent rewards:\par
Focal action | Partner cooperates | Partner defects\par
Cooperate | 8 | 0\par
Defect | 10 | 5\par
For any forecast probability q that the partner cooperates, if that forecast is independent of the candidate focal action:\par
- Expected cooperation reward is `8q`.\par
- Expected defection reward is `10q + 5(1-q) = 5 + 5q`.\par
- Defection’s advantage is `5 - 3q >= 2` for every q between 0 and 1.\par
Thus a stagewise rational response always defects, regardless of whether the next-action forecast is correct. Against the actual fixed-action partners, always defecting is also optimal over the whole episode. Under this specific reading, prediction-based regret must be zero. Sampling error or a confidence interval cannot explain a nonzero value for this exact deterministic identity.\par
\subsubsection*{Why that is not an unconditional proof of a bug}
The general notation allows a model of an entire counterfactual partner policy, not only a forecast independent of the focal action. An incorrect model could predict retaliation after defection, making cooperation rational under the model even when the true partner is fixed.\par
An independent 100-step example makes this distinction concrete. Let the true partner always cooperate. Let the imagined partner cooperate until the focal agent first defects, then defect forever from the next round. The optimal response to the imagined policy is to cooperate for 99 rounds and defect in the last round. Its reward against the true partner is `99*8 + 10 = 802`; the true optimal always-defect policy earns 1,000. Regret is therefore 1.98 per step, despite correct partner predictions along the chosen path. A small independent backward-induction calculation verifies the policy and rewards [E8].\par
This example satisfies the broader full-policy interpretation and defeats the claim that every possible rational model-based response must have zero regret. It is a diagnostic counterexample, not an explanation of the published numbers.\par
\subsubsection*{Supporting checks}
The prompt images elicit the next partner action. Figure 2 includes the focal agent’s already selected current action, while Figure 3 predicts the partner first; Figure 18 passes a single predicted action into an action-selection prompt. These differences matter. They do not by themselves document a counterfactual rollout model, how candidate actions are queried, a planning horizon, or a transition/value computation for the reported comparator.\par
The preserved arXiv source archive contains manuscript, bibliography/style material and figure assets, with no experiment implementation located. Its active metric text supplies no additional estimator or planner specification. A publisher/paper/author-page search did not locate a clearly linked implementation; OpenReview was blocked by a verification wall. These are bounded search results, not a claim that code does not exist.\par
Required resolution. Supply the exact estimator or relevant code excerpt, the information supplied to its optimizer, treatment of counterfactual actions and continuation values, and an example episode connecting forecasts to the reported regret. If the comparator is stagewise, correct the single-action IPD values to the appropriate zero-regret baseline and recompute dependent summaries. If it is a full-policy planner, document and validate that procedure. Until then, the numerical comparator and inferences specifically relying on it remain unverified.\par
This question is central enough to require resolution. It is not labeled a demonstrated Major because the alternative semantics are real and the historical calculation has not been identified.\par
\subsection*{6. Additional checks and interpretation limits}
Task relevance. The p. 4 relevance metric compares optimal responses across partners. Under a uniform distribution over the three fixed RPS actions, its expected per-step value is 1, matching the paper. For fixed-action IPD partners it is 0, because the same always-defect policy is optimal for both. Static IPD is consequently a useful control for payoff-directed decision making, but provides no partner-dependent optimal-action distinction under that metric. Dynamic tit-for-tat partners require a separate analysis; the zero result is not transferred to them.\par
Prediction versus action. A good next-action predictor need not be a good controller, and the paper explicitly permits model-free policies. A constant-action heuristic therefore does not refute its stipulated literal-ToM definition. Likewise, poor oracle-assisted choices can expose payoff reasoning or context-handling problems. They do not, by themselves, isolate a failure to infer another agent’s hidden mental state. The paper’s functional objective intentionally includes decision competence.\par
Reward versus regret. At a fixed task and comparator, cumulative regret is the optimal return minus the achieved return. The p. 4 wording about a constant factor should be expressed as a sign reversal plus an additive task-specific constant. Rankings correspond when the comparator is fixed; pooling tasks with different comparators needs care. No additional demonstrated table error follows from this wording alone.\par
Scope of the headline. Section 5, p. 9, explicitly allows insight-only uses of literal ToM and defends its stronger terminology for applications in which AI receives agency. Section 3 also acknowledges prior interactive and behavioral game-theory work. These qualifications constrain the paper’s claim and must remain in a faithful review. The recommendation for direct action evaluation is not refuted by the two local corrections.\par
Experimental coverage. Appendix C specifies 100-step episodes, a 24-hour or 100-episode limit and at least 30 episodes per model. The review read reported aggregates and selected curves, including Figures 19, 28-29 and 42-43. It did not reproduce model outputs, reconstruct confidence intervals, validate the RMax implementation or recover exact per-condition episode counts. Figure trends over 100 steps should be described at that finite horizon; an asymptotic convergence conclusion would need separate support. The reasoning-model successes reported in Table 4 also preclude an accurate summary that every model always fails.\par
\subsection*{7. Conclusion table}
Item | Supported judgment | Verification completed | Remaining action\par
Prior studies described as noninteractive | Minor | Three versioned primary protocols checked | Correct characterization and explain the added diagnostic\par
Figure 13 target mismatch | Minor in printed protocol | Original prompt image checked | Verify actual run template; correct figure or results as warranted\par
Nonzero single-action IPD prediction-based regret | Needs Evidence | Stagewise zero result and full-policy counterexample independently checked | Identify and validate the actual estimator\par
RPS task relevance | No issue found under uniform sampling | Exact finite enumeration gives 1 | State the sampling distribution\par
Broader functional evaluation recommendation | Not refuted by these checks | Scope and alternative views assessed | Human scientific assessment\par
Experiments and confidence intervals | Not independently reproduced | Reporting and selected figures inspected | Run-level evidence and implementation remain unverified\par
No author correspondence, paper modification or empirical repair was performed. The report separates demonstrated defects from conditional diagnoses and open evidentiary questions. Human scientific decisions, author adoption and publication status remain pending.\par
\subsection*{8. References and retained evidence}
- E1: Exact target `papers/0060/0060.pdf`; source-acquisition record, `identity\_verification.json`, extracted text and original-page renders.\par
- E2: https://proceedings.mlr.press/v267/riemer25a.html ; preserved publisher snapshot.\par
- E3: https://arxiv.org/abs/2412.19726 ; separately saved v4 PDF, 12 June 2025. SHA-256: `9d469d20a75ed81a4727079c27c932c107ba93db9e12094ab5767d6987f07e2b`. Targeted comparison: `verification/0060/arxiv\_v4\_targeted\_comparison.json`.\par
- E4: arXiv v4 TeX source, https://arxiv.org/src/2412.19726v4 ; archive SHA-256: `51d05737e9ab73fbf1ea801769d77dd93e10a2c75133e5b7f58792c950c2b3be`. Static evidence and search limits: `verification/0060/technical\_audit/`.\par
- E5: Xie et al., Can Large Language Model Agents Simulate Human Trust Behaviors?, https://arxiv.org/abs/2402.04559v1 ; PDF SHA-256: `d08e99db82a95e8807759d9813cdaa2fa58068b2000fe856bae8bd83fa931b70`.\par
- E6: Ha et al., CloChat, https://arxiv.org/abs/2402.15265v1 ; https://doi.org/10.1145/3613904.3642472 . PDF SHA-256: `19ad72fd6faa8c467d991b5a9fc5c4d9fc88eca6a4b34f6c3f2185469768f6e5`.\par
- E7: Li et al., Large Language Model-Empowered Agents for Simulating Macroeconomic Activities, https://arxiv.org/abs/2310.10436v1 . PDF SHA-256: `efcd6f8233b2520578db506b57b8d21a2a1fc2a51a29e042568ad57a35539ea4`. E5-E7 PDFs, receipts and exact locators are retained under the argument evidence directories.\par
- E8: Reviewer-written `verification/0060/independent\_metric\_checks.py` and its JSON output; exact payoff checks, finite task-relevance enumeration and backward induction. These are independent calculations, not execution of an author implementation.\par
- Workflow: Paper 01 / Paper 02 evidence guidance; Kassis et al. (2026), Scientific Agent Skills, https://doi.org/10.48550/arXiv.2609.00065 ; procedural reference only.\par
\subsection*{Supporting checks}
Original and separately acquired source bytes were retained with SHA-256 receipts. Human scientific assessment remains pending.\par
No author code was run or imported.\par
\end{document}
